% =====================================================================
%  Dichiarazione sull'uso di strumenti di IA generativa
%  Inclusa da main.tex subito dopo l'indice.
%
%  Titolo di sezione e non di capitolo, e nessun \clearpage: così la
%  dichiarazione prosegue sulla stessa pagina in cui l'indice finisce,
%  invece di aprirne una nuova quasi vuota. Senza voce nell'indice: è una
%  dichiarazione, non una parte del lavoro.
% =====================================================================

\bigskip
\noindent\rule{\textwidth}{0.4pt}

\section*{Dichiarazione sull'uso di strumenti di intelligenza artificiale
generativa}

Nella preparazione di questo lavoro è stato impiegato un modello linguistico di
grandi dimensioni, Claude Opus 5 di Anthropic, come strumento di supporto.

L'impiego ha riguardato due ambiti. Il primo è la scrittura e il collaudo del
codice di analisi e di generazione, a partire da specifiche formulate
dall'autore. Il secondo è la stesura e la revisione del testo, condotte su
contenuti, risultati e argomentazioni già stabiliti dall'autore.

Restano interamente dell'autore l'ideazione del lavoro, la formulazione della
domanda di ricerca, il disegno sperimentale, la ricerca bibliografica,
l'esecuzione delle misure, l'interpretazione dei risultati e le conclusioni che
ne discendono. Nessuna fonte è stata reperita, selezionata o citata per
suggerimento dello strumento, e nessun risultato numerico riportato in questa
tesi proviene da esso: tutti i valori sono prodotti dalle procedure descritte
nel Capitolo~\ref{cap:metodi} ed eseguite sui dati elencati
nell'Appendice~\ref{app:codice}.

L'autore si assume la piena responsabilità dell'intero contenuto della tesi.
